% Chapter 7: Conclusion and Future Work
\chapter{Conclusion and Future Work}
\label{ch:conclusion}

% ─────────────────────────────────────────────────────────────
\section{Summary of Achievements}
\label{sec:conc-summary}

This project set out to design, implement, and evaluate OrderVista, an
AI-enhanced logistics dispatch platform that unifies administrative task
management and driver delivery execution within a single Flutter application
backed by Supabase and DeepSeek~R1.

All fourteen objectives defined in Section~\ref{sec:objectives} have been
achieved and verified with traceable test evidence
(Table~\ref{tab:obj-map}).  The key outcomes are summarised below:

\begin{enumerate}
  \item A \textbf{unified, role-gated platform} serving both dispatchers and
        drivers from a single codebase, eliminating the need for separate
        applications.
  \item A \textbf{nearest-neighbour route optimisation engine} that reduces
        total delivery distance by 25--31\% compared to unoptimised task
        ordering (Table~\ref{tab:route-eval}), with sub-second on-device
        execution enabling mid-shift re-sequencing.
  \item An \textbf{agentic AI assistant} powered by DeepSeek~R1 that
        provides natural-language task summaries, contextual chat, and
        structured action recommendations with a human confirmation gate,
        achieving 6/6 correct action selections in evaluation
        (Table~\ref{tab:agentic-eval}).
  \item A \textbf{multi-tenant security model} enforced by 17 PostgreSQL
        Row-Level Security policies across 11 tables, verified through
        cross-tenant isolation tests (Section~\ref{sec:rls-test}).
  \item An \textbf{offline-first architecture} using Drift (SQLite) with
        exponential-backoff synchronisation, ensuring drivers can complete
        deliveries without network connectivity.
  \item \textbf{Real-time communication} between admins and drivers via
        Supabase Realtime WebSocket channels, including direct messaging
        and broadcast capabilities.
  \item A \textbf{driver recommendation system} that ranks available
        drivers by a composite score incorporating proximity, availability
        status, and historical delivery performance
        (Equation~\ref{eq:rec-score}).
\end{enumerate}

% ─────────────────────────────────────────────────────────────
\section{Contributions}
\label{sec:conc-contributions}

OrderVista makes three contributions to the intersection of mobile
software engineering and logistics technology:

\begin{enumerate}
  \item \textbf{Agentic AI in last-mile logistics}: to the author's
        knowledge, OrderVista is the first open-source logistics
        platform to integrate an agentic AI layer that can recommend
        \emph{and execute} operational actions (e.g.\ status transitions,
        sync triggers) subject to human confirmation.  This goes beyond
        the read-only summarisation offered by existing commercial tools.
  \item \textbf{Database-enforced multi-tenancy for fleet platforms}: the RLS-based isolation model demonstrates that robust data
        partitioning can be achieved without a custom middleware layer,
        providing a reusable pattern for BaaS-backed SaaS applications.
  \item \textbf{Offline-capable delivery completion with deferred sync}: the Drift queue with exponential backoff provides a reference
        implementation for Flutter applications that must tolerate
        intermittent connectivity in urban environments.
\end{enumerate}

% ─────────────────────────────────────────────────────────────
\section{Future Work}
\label{sec:future-work}

Several extensions would strengthen OrderVista for production deployment:

\paragraph{Short-term improvements.}
\begin{itemize}
  \item \textbf{2-opt post-processing}: adding a local-search pass
        after the nearest-neighbour heuristic could yield a further
        5--10\% distance reduction with minimal latency impact
        (Section~\ref{sec:disc-routing}).
  \item \textbf{Push notifications}: integrating Firebase Cloud
        Messaging (FCM) to alert drivers of new task assignments and
        admin messages when the app is in the background.
  \item \textbf{Server-side workflow validation}: implementing
        PostgreSQL \texttt{BEFORE UPDATE} triggers to enforce valid
        state transitions (e.g.\ preventing \texttt{delivered} without
        a PoD URL), closing the gap identified in
        Section~\ref{sec:disc-baas}.
  \item \textbf{Stale-cache warning}: displaying a time-since-last-sync
        banner in the driver app to mitigate the risk of acting on
        outdated task data during extended offline periods.
\end{itemize}

\paragraph{Medium-term extensions.}
\begin{itemize}
  \item \textbf{Time-window constraints}: extending the route model to
        respect customer-specified delivery windows, transforming the
        problem from a pure TSP to a TSP with Time Windows (TSPTW).
  \item \textbf{ETA calculation}: leveraging Mapbox Directions API to
        provide real-time estimated arrival times visible to both
        dispatchers and customers.
  \item \textbf{iOS deployment}: building and testing the iOS target
        with an Apple Developer account to validate cross-platform
        parity.
  \item \textbf{Google Play Store release}: a signed Android App Bundle
        has been uploaded to Google Play Console and a closed testing
        track with 12 recruited testers is currently in progress.
        Google's policy requires new developer accounts to maintain a
        closed test for a minimum of 14~days before production access
        is granted.  Full public release, including privacy-policy
        hosting and data-safety declarations, will follow upon
        completion of this mandatory review period.
\end{itemize}

\paragraph{Long-term research directions.}
\begin{itemize}
  \item \textbf{On-device SLM inference}: as small language models
        improve in structured-output reliability, migrating the agentic
        layer to on-device inference would eliminate cloud dependency
        and reduce per-query cost to zero.
  \item \textbf{Reinforcement learning for dispatch}: training an RL
        agent on historical delivery data to learn assignment policies
        that minimise fleet-wide completion time, surpassing the
        heuristic composite score currently used.
  \item \textbf{Multi-vehicle routing}: upgrading from single-driver
        TSP to full VRP with capacity constraints, enabling the
        optimiser to jointly schedule multiple drivers across a shared
        task pool.
\end{itemize}

% ─────────────────────────────────────────────────────────────
\section{Closing Remarks}
\label{sec:closing}

OrderVista demonstrates that a single-developer academic project can
deliver a feature-rich, AI-enhanced logistics platform by leveraging modern
BaaS infrastructure, open-weight language models, and disciplined iterative
development.  The platform's 14 fulfilled objectives, 29 passing functional
tests, and measurable route-optimisation improvements provide evidence that
the approach is sound.  While production readiness requires the
enhancements outlined above, the architectural foundations (modular route
pipeline, database-enforced tenancy, and human-gated agentic AI) offer a
solid base for continued development.

Ultimately, the project's value extends beyond its technical contributions.
By consolidating fragmented workflows into a single interface, providing
optimised routing and clear task sequencing, and placing AI-driven decision
support directly in drivers' hands, OrderVista aims to reduce the
avoidable coordination burden borne by frontline logistics workers, helping
them complete delivery work more safely, more efficiently, and with less
operational friction.
